\section{Terminal interpretation and independently replayed certificates}
\label{sec:terminal-certification}

\subsection{What the graph query means topologically}

The graph query used by the prototypes asks whether a cubic graph has a cut
of size at most three with an edge on each shore. Loops count twice toward
degree and never cross a cut; parallel edges remain distinct occurrences.
For a cubic graph this is exactly a cut with a cycle on each shore. Indeed,
if a shore induces a forest with $v$ vertices and $a$ components, then
\[
 |\delta X|=3v-2(v-a)=v+2a.
\]
A forest containing an edge has $v\geq2$, $a\geq1$, and hence boundary at
least four. The argument explains why the two monochromatic-edge marks carry
the needed information at threshold three. At threshold four they do not
have the same interpretation.

For an actual cubic pattern $P\subset\partial B^3=S^2$, essentiality means
that every transverse simple closed curve crossing at most three pattern
edges bounds, on at least one side, a disc whose pattern is empty, an arc,
or a tripod. The incoming primal criterion states that essentiality is
equivalent to absence of the marked cut just described
\cite{incoming-four-copy,lackenby2026}.

For completeness, a violating curve has a cycle on either side: a forest in
a disc with at most three boundary leaves can only be empty, an arc, or a
tripod, since a tree with $v$ trivalent interior vertices has $v+2$ boundary
leaves. Coloring graph vertices by side of the curve gives a marked cut of
size at most three. Conversely, in a connected plane graph choose a smallest
cut having cycles on both shores. Both shores are connected; otherwise the
boundary of a cycle-containing component is smaller and the opposite cycle
persists. Such a cut is a bond. The boundary of a regular neighborhood of
one shore, with complementary regions away from the other shore filled,
gives a simple curve crossing exactly its edges. If the cubic graph is
disconnected, a regular-neighborhood boundary separates two cycle-containing
components at cost zero. This proves the criterion in all represented cases.

Ordinary edge connectivity would give the wrong predicate. The three edges
incident to a vertex of $K_4$ cut off an allowed tripod. A triangular prism
also has edge connectivity three, but its three rungs separate two cycles.
The marked query distinguishes them. A geometric importer must additionally
handle vertex-free circle components and the empty pattern; the closed-chain
API here encodes the nonempty cubic part.

\subsection{A finite graph certificate engine}

On a finite kernel of size $K$, enumerate pairs of edge occurrences with
disjoint endpoint sets. Require one edge to be monochromatic zero and the
other monochromatic one. Give the arcs fixing their endpoints capacity four,
and replace every non-loop undirected edge by two opposed unit-capacity arcs.
An integral flow computation need only find four units of flow or a smaller
minimum cut. In the latter case the induced coloring has the required marks.
If every anchor pair supports four units, no marked cut of size at most three
exists. This is the unchanged incoming finite certificate engine
\cite{incoming-four-copy}.

A negative graph proof lists a shore; replay recounts its exact cut and its
two edge marks. A positive graph proof supplies residual augmenting paths
reaching four for every disjoint anchor pair. The checker verifies those paths
and requires the complete set of pairs. It does not perform max-flow search.
With sparse flow construction, the reference producer has a conservative
$O(K^3)$ structural bound and $O(K^3\log(K+1))$ certificate size. The
independent dense reference replay has an $O(K^4)$ bound. These are adequate
polynomial bounds, not claims to optimal algorithms for cyclic connectivity.

In a source chain, private vertices must have degree three. At every serial
seam the corresponding output and input degrees must sum to three. The same
checks apply within every repeated block and at the final closure. Interfaces
are injective and disjoint before the stated seam quotient. These local
conditions prove that the actual expanded graph is cubic, including possible
loop edges after closure. The reblocking step is an exact rebracketing of that
quotient; the routing step retains the degrees at its genuine labeled seams.

\subsection{Three replay layers for the pattern kernel}

The pattern producer and checker share immutable input types and elementary
validation. They use separate component reconstruction and gluing procedures.
The producer follows component chains; replay independently identifies
components with a disjoint-set structure and assembles fresh-copy quotient
vertices. It verifies the finite graph proof on the reconstructed kernel.
For a negative result it also verifies local component color masks, consecutive
seams, repeated monochromatic cycle words, the exact binary count of original
copies, and the complete list of cut edges and monochromatic anchors.

The through-component wrapper adds a preceding exact reblocking layer. Its
certificate binds the original chain, a bounded aligned power of each tile,
the unchanged remainder, the adapted-source identity, and the nested routing
proof. Its checker constructs the bounded power independently and checks the
actual component partitions. The affine address maps in
Section~\ref{sec:through-reblocking} then recover original block, copy, and
edge-occurrence identifiers. Thus a cut on the reduced graph is accompanied
by a succinct coloring of the original graph, even when its exponent has
thousands of bits.

It is useful to keep the three correctness assertions separate:
\begin{enumerate}
\item The original tile grammar and the adapted grammar describe the same
finite graph before reduction, with an exact map of edges.
\item The adapted grammar and the bounded kernel have the same marked query
signature; source colors for a low cut lift by zero-cost repetition.
\item The finite graph proof establishes the query answer on that kernel.
\end{enumerate}
For a topological conclusion there is a fourth, geometric assertion: the
source grammar is the actual spherical pattern on an independently certified
ball. No graph certificate supplies it.

\subsection{Return values and the native integration contract}

The three new operators have intentionally different meanings.
\begin{center}\small
\begin{tabularx}{\linewidth}{@{}P{0.23\linewidth}P{0.27\linewidth}Y@{}}
\toprule
Operator & Successful output & Meaning of other outcomes\\
\midrule
Cancellation search & \texttt{DISC\_FOUND}, with source-bound normal
coordinates and vertex-disc proof & Exhaustion only covers the stated class;
limits and unsupported sources remain inconclusive\\
Bundle classification & \texttt{COMPLETE}, with exact component/genus census
and source proofs & Saturated bases or unresolved cores remain inconclusive\\
Pattern query & \texttt{LOW\_CYCLIC\_CUT} or
\texttt{NO\_LOW\_CYCLIC\_CUT}, with proof & Malformed or unsupported
large tiles are rejected; neither result alone is a knot verdict\\
\bottomrule
\end{tabularx}\end{center}

Discovery limits are charged globally within one query, including attempted
systems and profile combinations. External cancellation propagates. Replaying
a supplied positive certificate does not call the candidate search. Geometry
replay checks the underlying source-parallelity, boundary-word, and raw-core
proofs before computing the finite census. These choices make every positive
claim reviewable without treating an implementation timeout as evidence of a
negative topological conclusion.

The proposed repository integration is additive and opt-in. First import the
specific incoming dependencies, then add the new cluster and bundle modules,
then expose the graph operator at a boundary where its grammar contract is
proved. A native recognizer benchmark would be meaningful only after that
producer exists. The present measurements report the actual local operations
performed; they do not substitute synthetic graph times for knot-recognition
times.
